\section{A fourteen-coordinate equation of degree 56}
\label{sec:fourteen}

The quotient absorber also applies before radical elimination. Combined
with the inherited ordinary strong norm, it gives fourteen strictly
positive unknowns, exact degree~56, and 79 binary arithmetic operations.
The historical $Q_{56}$ has fifteen positive unknowns, degree~56, and
74 operations. Thus the construction exchanges five operations for one
fewer unknown, while retaining the same degree. This is a separate
tradeoff from the equations in eleven or fewer coordinates.

Keep the common definitions of $q,C,X,Y,E,k,c,a,H,\Delta,R,b_0$,
with $Y=sq^3$ and $b_0=X+ac$. The supplied coordinates are
\[
 J,F,\alpha,z,f,i,T,s,w,\tau,\eta,\zeta,y,\gamma>0.
\]
In addition to $U=4q^4$ and $r=k-R-1$, put
\begin{align*}
 V&=c(Tf-1)-Rf^2, & G&=1+\Delta c^2,\\
 N_0&=\tau^2-(XY^2)(XY^2+1)k^2,
 &I&=ic^2,\\
 A_{\rm aux}&=\Delta(f^2-1),
 &N_a&=A_{\rm aux}(V^2-y^2)+y^2,\\
 N_t&=(K_0+w)C+q-F-z(q-1),
 &D_*&=Ub_0+H(U\gamma-r).
\end{align*}
The single ordinary integer polynomial equation is
\begin{equation}
 \boxed{(D_*^2-U^2G)^2+(I^2-A_{\rm aux})^2+(N_a-1)^2+(N_0N_t-1)^2=0.}
 \label{eq:fourteen}
\end{equation}
Every abbreviation is computed, not another unknown. The same complete
package-31 coefficient tuple supplies all four fixed ports. The coordinate
$i$ now has the ordinary strong meaning: the equation is
$I^2=\Delta(f^2-1)$. One must use the ordinary strong-rank theorem
instead of canceling $\Delta$ informally.

\paragraph{An early lower bound on the index.}
At a zero of \eqref{eq:fourteen}, all four residuals vanish. The
consecutive-product negative-Pell descent excludes $N_0=-1$, so
$N_0N_t=1$ gives $N_0=N_t=1$. The authentic transport interface yields
$q\ge16$, $3q<R<q^4-2$, and $E\ge q^4>R+2$.
The main residual makes $g=D_*/U$ a rational square root of the
integer $G$, hence an integer. It is positive: $H<5a$, $b_0>akY$,
and $UY>5$ give $D_*>Ub_0-Hk>ak(UY-5)>0$.
The carrier implies $U\mid Hr$.

Write $A_*=a+2$ and $P_0=2XY^2+1$. Integer-coordinate Pell
classification supplies positive indices $p,n$ with
\[
 c=\psi_{A_*}(p),\quad g=\chi_{A_*}(p),\quad
 k=2\psi_{P_0}(n).
\]
Because $P_0>A_*$ and $c>k$, we have $p>n$.
The recurrence gives $k\equiv2n\pmod{4XY^2}$, and
$U\mid4XY^2$. Also $\gcd(U,H)\le3$: $H$ is odd, every
common prime divisor of $q$ and $H$ divides~3, and $3\mid q$
implies $H\equiv3\pmod{27}$. Therefore
\[
 \mu=U/\gcd(U,H)\ge4q^4/3>R+1.
\]
The divisibility $U\mid H(k-R-1)$ gives
$2n\equiv R+1\pmod\mu$. Since $n>0$ and $0<R+1<\mu$,
\[
 2n=R+1+j\mu\quad(j\ge0).
\]
Thus $n\ge25$ and $p\ge26$ \emph{before} recovering $p=R$.
The elementary recurrence bound
$c\ge(2A_*-1)^{p-1}\ge A_*^{25}$ gives
\[
 c>A_*\Delta^2,\qquad c>2p,\qquad c>2R.
\]
For the last inequality one may also use $c\ge2A_*>2R$.

\paragraph{Ordinary strong-rank recovery.}
Apply the inherited ordinary strong-rank lemma, with its strict size
hypothesis now proved. From
$c=\psi_{A_*}(p)>A_*\Delta^2$ and
$(ic^2)^2=\Delta(f^2-1)$ it gives a positive index $m$ such that
\[
 f=\chi_{A_*}(m),\quad p\mid m,\quad c\mid m,\quad
 m\ge c>2p,\quad I=\Delta\psi_{A_*}(m),\quad c^2\mid f^2-1.
\]
The proof, reproduced among the inherited sources, handles the possible
nonfundamentality of $A_*$ in the squarefree Pell group. Its large-$c$
condition is essential to this invocation. No assertion that $I/\Delta$
is integral follows merely from the original displayed strong equation.

Set $Q=I^2=A_{\rm aux}$. The auxiliary equation becomes
$QV^2-(Q-1)y^2=1$. Its odd Pell index $j_0$ satisfies
$j_0\equiv\pm p\pmod m$ by reduction modulo $f$ and signed
abscissa reduction; it satisfies $j_0\equiv\pm R\pmod c$ by
reduction modulo $c$. The latter uses the recovered divisibility
$c^2\mid f^2-1$. Since $c\mid m$ and $0<p,R<c/2$, it follows
that $p=R$. These are exactly the ordinary signed auxiliary steps in
the inherited compiler proof, with all required bounds available.

Now $1\le n<p=R$, so
$|2n-R-1|<R<\mu$. The earlier congruence forces $j=0$ and
$2n=R+1$. Hence $r=k-2n$ is divisible by both $E$ and $U$.
Define
\[
 h=r/E,\qquad \gamma_{\rm old}=\gamma-r/U=(g-b_0)/H.
\]
Both are integers. The first is positive since $n\ge25$ and
$\psi_{P_0}(n)>n$. Also $k\ge4P_0>X$ gives $c>X$, while
$g>(a+1)c$ implies $g-b_0>c-X>0$, proving
$\gamma_{\rm old}>0$. All fifteen positive coordinates of the
ordinary-strong blank compiler have been restored, so its calibrated
soundness theorem supplies the ZFC refutation.

\paragraph{Constructive converse.}
A positive ordinary-strong compiler witness satisfies $R=2n-1$.
Thus $r=k-2n$ is a positive multiple of $4XY^2$, hence of $U$.
Set $\gamma=\gamma_{\rm old}+r/U$ and retain the other thirteen
coordinates. Then $D_*=Ug$, and all four residuals in
\eqref{eq:fourteen} vanish. No new condition on $s$ or increase of a
Pell index is needed. Canonical normalized witnesses give ordinary
witnesses by the established substitution
$i_{\rm ordinary}=\Delta i_{\rm normalized}$.
This proves equivalent positive solvability and hence the stated
equivalence of unsolvability with $\operatorname{Con}(\mathrm{ZFC})$.

\paragraph{Exact degree and reproducibility.}
The four residual degrees are $26,22,28,24$. For the first use
\[
 D_*^2-U^2G
 =U^2\bigl((b_0+H\gamma)^2-G\bigr)
  -2UHr(b_0+H\gamma)+H^2r^2.
\]
The three terms have degrees $26$, at most $25$, and at most $20$;
the ordinary main residual in parentheses has exact degree~18.
Only $(N_a-1)^2$ reaches degree~56, with highest homogeneous form
\[
 T^4f^8w^4s^8(\eta+\zeta)^4(\beta J)^{28}.
\]
Its $J^{28}\eta^4T^4f^8w^4s^8$ coefficient is
$\beta^{28}\ne0$. The independent standard-library checker expands
the complete main residual into 2,994 monomials with the coefficient
ports symbolic, proves its degree~26, and checks all 79 live gates
and the final nonzero highest form. This finite certificate verifies
the expression and degree. The unbounded semantic dependencies remain
the inherited ordinary compiler and machine calibration.
